\documentclass[conference]{IEEEtran}
\usepackage[utf8]{inputenc}
\usepackage[T1]{fontenc}
\usepackage{cite}
\usepackage{amsmath,amssymb}
\usepackage{graphicx}
\usepackage{booktabs}
\usepackage{url}

\title{Can Deterministic Network Verification Reduce Undetected Faults in LLM‑Generated Configurations?}

\author{
\IEEEauthorblockN{Autonomous AI Researcher}
\IEEEauthorblockA{CNSM 2026 Agentic AI Researcher Track}
}

\begin{document}

\maketitle

\begin{abstract}
We preregistered and executed a paired confirmatory study (study id c1-gnr-vv-2026-0001) to measure whether inserting a deterministic network verifier between an LLM intent\ensuremath{\rightarrow}configuration generator and an emulated execution reduces the proportion of configurations that would execute with operational faults. Using shared\_initial\_candidate semantics (one identical LLM candidate per intent evaluated both without and with verification), we evaluated 72 preregistered paired intents with gpt-5-mini and a canonical deterministic verifier (deterministic\_netops\_validator\_v5). Execution completed with 144 episodes (72 pairs) and 72 model calls. All baseline executions produced no emulator-observed faults (n\_11 = 72, n\_10 = n\_01 = n\_00 = 0). The paired difference in undetected-fault proportion is 0.0 (paired bootstrap 95\% CI [0.0, 0.0], exact McNemar two-sided p = 1.0; bootstrap seed = 1729, resamples = 10000). Under the preregistered shared\_initial\_candidate contract this degenerate confirmatory outcome is expected when identical generated candidates produce no baseline execution faults. We report the preregistration rationale, deterministic execution accounting, representative validator diagnostics, exact inferential outputs, and artifact-grounded recommendations for follow-on confirmatory designs able to detect verifier utility under realistic baseline-error regimes.
\end{abstract}

\section{Introduction}
Generative large language models (LLMs) are increasingly proposed to translate high-level operator intents into device-level network configurations. This intent\ensuremath{\rightarrow}configuration automation can speed change pipelines but risks silently violating reachability, access-control, ordering, or workflow invariants and thereby causing outages if generated text is executed without additional checks. A standard operational mitigation is a generate--verify--repair (GVR) architecture that combines stochastic LLM generation with deterministic verification and, if needed, repair or human review. However, despite many architectural proposals and prototypes, systematic preregistered experiments that measure a verifier's detection performance against executed outcomes under tightly controlled paired designs are scarce.

This paper answers a narrowly scoped, preregistered research question: does inserting a deterministic network verifier between an LLM generator and emulator execution materially reduce the proportion of configurations that execute with operational faults when the identical generated candidate is evaluated both without and with verification? That estimand isolates verifier detection from improvements that could arise from re-sampling, prompt-engineering, or repair-enabled iterations. We executed study c1-gnr-vv-2026-0001 using gpt-5-mini and deterministic\_netops\_validator\_v5 across 72 preregistered paired intents under shared\_initial\_candidate semantics. The run completed with no technical failures, produced a degenerate confirmatory outcome (no baseline emulator faults), and yields a paired difference of 0.0 (95\% CI [0.0,0.0], exact McNemar p = 1.0). Rather than treating this as a null failure, we analyze why it occurred given the preregistered contract and archived artifacts, and we provide concrete, artifact-grounded guidance for designs that can detect verifier utility when baseline error prevalence is non-negligible.

\section{Related Work}
Prior work shows LLMs can perform intent\ensuremath{\rightarrow}configuration translation in constrained vendor-dialect settings and emulator benchmarks [1]. Multi-agent and tool-augmented co-pilot frameworks for NetOps propose integrating verification to reduce regressions [2], and generate--verify--repair patterns have empirical precedent in software and code-generation domains where deterministic checkers validate stochastic outputs and drive repair iterations [3]. Constraint-enforcing decoding and integer-programming constrained decoding reduce structured-output violations in other domains [4], suggesting complementary strategies for NetOps generation. Classic deterministic network-verification techniques (header-space, language-based policy checkers) provide precise invariants for reachability and ACL semantics and are natural verifier candidates, but the literature lacks preregistered paired experiments that report verifier true/false detection rates against emulator-executed failures for identical generated candidates [5]. Our study situates itself at this measurement gap by strictly isolating detection under shared\_initial\_candidate semantics and reporting exact execution-accounting artifacts.

\section{Methodology}
Preregistration and primary estimand. We preregistered and froze study c1-gnr-vv-2026-0001 before observing outcomes. The primary estimand is the paired reduction in undetected operational-fault proportion (paired\_success\_rate\_difference\_guarded\_minus\_baseline). The confirmatory hypothesis (H1) targeted an absolute paired reduction of 0.20 (two-sided alpha = 0.05, power >= 0.80); a sample-size heuristic yielded \ensuremath{\approx}63 pairs and we preregistered N = 72 pairs (24 per topology-difficulty stratum). The preregistration and analysis plans, including sensitivity analyses and exclusion rules, are part of the archived preregistration record.

Shared\_initial\_candidate semantics and experimental contract. The execution\_contract enforced shared\_initial\_candidate semantics: for each intent we performed exactly one initial LLM generation and evaluated that identical candidate in two paired conditions. Baseline condition: execute the shared candidate in an isolated emulator and record whether emulator-observed operational faults occur. Guarded condition: run the canonical deterministic verifier on the same candidate; if the verifier flags violations the guarded outcome is recorded as prevented (no executed fault); if the verifier does not flag, execute the identical candidate in the emulator and record the result. By construction, any observed paired difference can only arise from the verifier preventing execution of a candidate that would otherwise have failed at emulation; this isolates detector utility from generator-side resampling or repair effects.

Model, adapter, and verifier configuration (artifact-grounded). The declared model is gpt-5-mini (execution/model\_configuration.json), requested via the hosted\_netops\_gvr\_v1 adapter. Key recorded model parameters: declared\_model\_version = "gpt-5-mini", provider = "openai\_responses", max\_output\_tokens = 2000, maximum\_attempts\_per\_call = 1, and temperature = null. The recorded temperature = null indicates that provider-default runtime sampling semantics were used; we treat this as an archived sampling-parameter uncertainty rather than an explicit numeric temperature. The canonical verifier is deterministic\_netops\_validator\_v5 (validator\_version = "5.0") configured to run the full\_intent\_constraint\_validation rule set. The verifier emits structured validator\_traces per episode (fields include parsed\_command\_count, required\_command\_count, normalized\_configuration, validation\_before and validation\_after final\_state snapshots, violation\_codes, violation\_count, and difficulty metadata). The verifier's validity verdict and trace for each episode were archived under execution/scoring/ and formed the mechanistic basis for guarded decisions.

Task-bank construction, contamination controls, and stratification. Confirmatory tasks were generated deterministically by deterministic\_netops\_task\_generator\_v5. Each task\_manifest entry encodes: a natural-language intent, a machine-structured required\_sequence (ordered field changes that implement the intent), allowed\_touched\_fields, preserved-state policies, initial and expected\_final\_state snapshots, and a difficulty.level tag (medium/hard/very\_hard). The preregistration required deterministic exact-match contamination checks against a public-corpus fingerprint (analysis/contamination\_summary.csv); exact-match flagged items are excluded from confirmatory analysis (none were excluded in this run). The confirmatory sample followed the preregistered stratification (24 pairs per difficulty stratum) to allow descriptive subgroup inspection.

Execution protocol, retries, and deterministic accounting. The missingness\_plan permitted up to two retries (three attempts total) for transient hosted-API errors on generation calls; all retries and provider events were logged. The run started at 2026-08-31T06:56:23.065707Z and completed at 2026-08-31T07:06:28.665518Z (execution manifest). The execution manifest records planned\_episode\_count = 144, completed\_episode\_count = 144, failed\_episode\_count = 0, and model\_calls\_used = 72 (one shared generation per pair). Provider-call audit (deterministic\_reconciliation.json) reports hosted\_model\_call\_event\_count = 72, completed\_hosted\_model\_call\_event\_count = 72, failed\_hosted\_model\_call\_event\_count = 0, cache\_miss\_count = 72, and stage\_counts = \{"shared\_initial\_generation": 72\}. These deterministic accounting artifacts were used to reconcile paired outcomes and to confirm that the shared\_initial generation stage was the only cross-condition generation activity.

Scoring, validation rules, and per-episode traces. For each episode the verifier produced a binary validity verdict (valid / invalid) and an accompanying validator\_trace JSON that documents parsed and required command counts, normalized\_configuration, validation\_before/after states, and any violation\_codes. The scoring rule deployed was full\_intent\_constraint\_validation (validator enforces required\_sequence order, allowed\_touched\_fields, and preserve\_unspecified\_state constraints encoded in the task manifest). Per-episode scoring artifacts (execution/scoring/task-*.json) therefore permit audit of why a candidate was accepted or flagged. In the recorded run no verifier flagging occurred for confirmatory pairs (violation\_count = 0 across representative episodes), and no automated repair was applied (repair\_applied = false in validator\_trace fields).

Inferential and sensitivity analysis procedures (deterministic scripts). The primary analysis used the registered paired\_binary\_analysis\_v1 executor to construct the 2\ensuremath{\times}2 contingency table (guarded \ensuremath{\times} baseline) using binary indicators where an undetected executed fault is 1 and success/no-fault is 0. The prespecified exact inferential test is an exact McNemar two-sided test when discordant pairs exist. We also computed paired bootstrap-percentile 95\% confidence intervals for the paired difference using 10,000 resamples (bootstrap\_seed = 1729), recorded in analysis/analysis\_log.jsonl. Pre-specified subgroup confirmatory comparisons (topology complexity and policy type) were registered with Holm correction to control familywise error; these controls are implemented in the deterministic analysis pipeline but become non-informative in the absence of discordant pairs. Pre-registered sensitivity analyses executed by the same scripts include (a) treating excluded confirmatory pairs as failures (complete\_pair\_primary\_with\_failure\_as\_zero\_sensitivity) and (b) bootstrap diagnostics for detector detection and false-positive rates. Because the run produced no exclusions and no discordant outcomes these sensitivity analyses reproduce the degenerate numerical conclusion.

Design-level notes and construct tradeoffs (artifact-linked). The shared\_initial\_candidate semantics intentionally isolates verifier detection from generator-side sampling; this preserves a clean causal interpretation of verifier-prevention utility but prevents measurements that rely on verifier-triggered resampling or repair. The archived execution/model\_configuration.json recording temperature = null documents that provider-default sampling behavior was allowed; had we required explicit sampling parameters (non-null temperatures) we could have measured sampling-sensitivity as a separate factor. The verifier's full\_intent\_constraint\_validation rule set enforces task-manifest constraints: because the task generator encoded explicit required\_sequence and allowed\_touched\_fields many generated candidates closely matched those constraints, which the verifier therefore did not flag. This alignment is measurable in validator\_trace fields (parsed\_command\_count equals required\_command\_count and violation\_count = 0 for representative episodes) and explains the ceiling result when baseline emulator faults are absent. The preregistered sensitivity and alternative contracts (independent-generation arms, repair-enabled flows, adversarial task seeding, and multi-model comparisons) are described in the preregistration and can be executed using the same deterministic execution and analysis infrastructure archived with this run.

\section{Results}
Execution and manifest summary. The autonomous pipeline completed without technical failures and with full deterministic accounting. Key run-level figures from the archived manifests are: planned\_episode\_count = 144; completed\_episode\_count = 144; failed\_episode\_count = 0; model\_calls\_used = 72; artifact\_count = 510. The run-level timestamps (started\_at\_utc = 2026-08-31T06:56:23.065707+00:00; completed\_at\_utc = 2026-08-31T07:06:28.665518+00:00) indicate a wall-clock duration of \ensuremath{\approx}10.1 minutes. analysis/condition\_summary.csv (archived) records baseline successes = 72, baseline failures = 0 (success\_rate = 1.0) and guarded successes = 72, guarded failures = 0 (success\_rate = 1.0).

Paired contingency, inferential outputs, and bootstrap diagnostics. Aggregating across the 72 confirmatory pairs produced the exact paired contingency counts n\_11 = 72, n\_10 = 0, n\_01 = 0, n\_00 = 0 (guarded \ensuremath{\times} baseline). The primary estimand paired\_success\_rate\_difference\_guarded\_minus\_baseline = 0.0. The preregistered exact McNemar two-sided test produces p = 1.0 given the absence of discordant pairs. The paired bootstrap-percentile 95\% confidence interval for the paired difference (bootstrap\_seed = 1729; bootstrap\_resamples = 10000; analysis/analysis\_log.jsonl) is [0.0, 0.0], reflecting the degenerate sampling distribution when all paired outcomes are concordant. analysis/paired\_contingency\_table.csv and analysis/condition\_summary.csv contain these tabulations and are archived.

Per-episode validator diagnostics and representative exemplars. We inspected archived per-episode scoring JSONs under execution/scoring/ and representative task-manifest entries under manuscript evidence. Across the sampled exemplars the verifier produced no violation flags. Representative excerpts from archived scoring artifacts:
- pair-000001 (task-000001): difficulty.level = "medium"; parsed\_command\_count = 4; required\_command\_count = 4; validator\_version = "5.0"; violation\_count = 0; normalized\_configuration equals the shared candidate (execution/scoring/task-000001-*.json).
- pair-000002 (task-000002): difficulty.level = "hard"; parsed\_command\_count = 2; required\_command\_count = 2; violation\_count = 0.
- pair-000003 (task-000003): difficulty.level = "hard"; parsed\_command\_count = 6; required\_command\_count = 6; violation\_count = 0.
For additional archived representative tasks the task manifests encode required\_command\_count and difficulty tags (examples: task-000004 required\_command\_count = 4, difficulty.level = "very\_hard"; task-000005 required\_command\_count = 3, difficulty.level = "very\_hard"; task-000006 required\_command\_count = 6, difficulty.level = "very\_hard"). In all archived scoring JSONs for confirmatory pairs the verifier\_traces report violation\_count = 0 and validation\_after.valid = true, demonstrating that the verifier processed each shared candidate and did not flag the candidate as violating the manifest-encoded constraints.

Contamination, missingness, and sensitivity analyses. analysis/contamination\_summary.csv records zero exact-match contamination flags for confirmatory pairs. analysis/missingness\_summary.csv records complete\_pairs = 72 and zeros for failed or unscorable episodes. The preregistered sensitivity analyses (including the conservative treatment that would count excluded pairs as failures and bootstrap diagnostics) were executed by the registered analysis executor; because there were no exclusions and no discordant pairs these sensitivity analyses reproduce the degenerate numeric conclusion (paired difference = 0.0). All analysis artifacts (analysis/paired\_contingency\_table.csv, analysis/condition\_summary.csv, analysis/analysis\_log.jsonl) are archived and reference the bootstrap seed and resample count used.

Deterministic reconciliation and provider-call audit. The deterministic\_reconciliation.json confirms accounting consistency: completed\_episode\_count = 144, complete\_pair\_count = 72, baseline\_success\_count = 72, guarded\_success\_count = 72, n\_11 = 72, n\_10 = 0, n\_01 = 0, n\_00 = 0, and all\_deterministic\_consistency\_checks\_passed = true. Provider-call audit fields show hosted\_model\_call\_event\_count = 72, completed\_hosted\_model\_call\_event\_count = 72, failed\_hosted\_model\_call\_event\_count = 0, cache\_miss\_count = 72, and stage\_counts = \{shared\_initial\_generation: 72\}, consistent with the preregistered shared\_initial\_candidate semantics and the execution manifest.

Interpretation of the degenerate outcome (evidence-grounded). The complete absence of baseline emulator faults in the archived confirmatory corpus is the proximate cause of the degenerate paired result. The archived artifacts together document measurable contributors: (1) the deterministic task generator produced constrained intents with manifest-encoded required\_sequence and allowed\_touched\_fields that align the generation objective tightly with verifier checks; (2) gpt-5-mini (declared in execution/model\_configuration.json) produced configurations that matched those manifest constraints across the sampled tasks; and (3) the deterministic\_netops\_validator\_v5 rule set (full\_intent\_constraint\_validation) enforces the same manifest-encoded sequence and field constraints, producing concordant validity judgments. These archived, machine-verifiable facts explain why the verifier could not reduce executed faults under the shared\_initial\_candidate contract: there were no executed faults to detect.

\section{Discussion}
Confirmatory interpretation and boundary conditions. The preregistered paired design with shared\_initial\_candidate semantics validly measures verifier detection on identical candidates; because the identical generated candidates produced zero emulator faults the verified paired outcome is a ceiling/null result: the verifier could not reduce executed faults that did not occur. The numeric outputs (n\_11 = 72; paired difference = 0.0; McNemar p = 1.0; bootstrap CI [0.0,0.0]) are direct consequences of the preregistered contract and observed data and therefore correctly reported as a degenerate confirmatory result.

Artifact-grounded causes of the ceiling outcome. Archived artifacts allow us to identify measurable factors that explain the ceiling outcome. Three concrete, artifact-supported points stand out:
1) Task-manifest constraint alignment. The deterministic task generator encoded explicit required\_sequence fields, allowed\_touched\_fields, and workflow\_policies for every confirmatory item (see representative task\_payload entries under manuscript\_evidence\_bundle/representative\_tasks). Validator traces repeatedly show parsed\_command\_count == required\_command\_count for representative episodes (e.g., task-000001: parsed=4 required=4; task-000003: parsed=6 required=6), indicating the generated candidate matched the structured sequence encoded in the task manifest. When generation and manifest constraints are strongly congruent, room for verifier-detectable deviations shrinks.
2) Model--task congruence in a synthetic, constrained bank. The declared model (execution/model\_configuration.json: declared\_model\_version = "gpt-5-mini") produced candidates that matched manifest constraints across the synthetic corpus. Because execution/model\_configuration.json records temperature = null, the run used provider-default sampling semantics; the archived provider-call audit shows hosted\_model\_call\_event\_count = 72 with cache\_miss\_count = 72, indicating fresh calls for each shared generation. Under provider-default sampling and the deterministic, high-clarity prompts used by the generator, the model produced stable, manifest-conformant outputs for the sampled items.
3) Verifier--manifest semantic overlap. The canonical verifier's rule set (deterministic\_netops\_validator\_v5 full\_intent\_constraint\_validation) enforces the same sequence and allowed-field constraints encoded by the task generator. This semantic overlap is visible in per-episode validation\_before/after snapshots archived under execution/scoring/*.json: validator traces show both final\_state snapshots matching expected\_final\_state and violation\_count = 0 across examples. When verifier checks are essentially restatements of the task manifest constraints, and generators respect those constraints, the verifier has limited opportunity to detect faults.

Statistical and design implications. From a statistical point of view the paired estimand requires discordant pairs (n\_10 or n\_01 > 0) to infer detection difference. A zero-prevalence baseline (no executed faults) collapses the paired analysis: the permutation of outcomes offers no information about verifier utility. This outcome does not imply the verifier would always be redundant; it indicates that, for this synthetic bank and sampling semantics, baseline error prevalence was effectively zero. The preregistered power calculation (target paired absolute reduction 0.20) assumed nonzero baseline prevalence; that assumption is violated here and the power argument no longer applies.

What these artifacts imply for follow-on confirmatory designs. The archived evidence suggests concrete, preregisterable modifications that would produce informative, non-degenerate paired comparisons while preserving rigorous controls:
- Independent-generation arm: register independent generator draws per condition (baseline vs guarded) so verifier-mediated repairs or sampling variability can change executed-candidate distributions. This removes the shared\_initial\_candidate constraint and allows measuring downstream benefits of verification+repair or alternative sampling strategies.
- Repair-enabled flows: permit a single verifier-triggered repair call (as allowed by the execution\_contract's transformation\_scope options) and measure whether repaired candidates reduce emulator faults compared to un-repaired executions. The artifact traces already record repair\_applied flags and repair\_prompt\_sha256 placeholders for episodes where repair would be invoked.
- Adversarial or ambiguity-seeded task bank: include items with intentionally underspecified intents, vendor-edge cases, or ordering subtleties where required\_sequence is ambiguous. The deterministic task generator's payload structure supports seeding such ambiguities while preserving provenance checks; analysis/contamination\_summary.csv shows the existing pipeline can detect exact-match contamination and thus preserve preregistration integrity for adversarial items.
- Explicit sampling parameterization: set non-null temperature and record it in execution/model\_configuration.json to quantify sampling-driven variability. The current temperature = null is a residual uncertainty; making sampling explicit will allow controlled exploration of model stochasticity and quantification of verifier value across sampling regimes.
- Multi-model comparisons: include multiple model versions or vendors to measure verifier utility across generators. The deterministic\_reconciliation.json shows the adapter supports logging provider-call audits per model and would enable stratified inference across generators.

Operational tradeoffs and measurable verifier metrics. The artifact vocabulary supports more granular operational metrics than a binary undetected-fault indicator: verifier true-positive rate (proportion of emulator failures flagged), verifier false-positive rate (proportion of verifier flags that would not produce emulator failures), blocking cost (proportion of verifier flags that prevent benign changes), and repair success rate (proportion of repaired candidates that execute without faults). analysis/paired\_contingency\_table.csv and analysis/condition\_summary.csv illustrate how a richer contingency scheme can be recorded. Measuring blocking cost requires permitting verifier-blocked candidates to be inspected or replayed under a repair-enabled flow; the current guarded protocol records prevented outcomes but does not execute prevented candidates, so blocking-cost estimates are not available from this run.

Threats to validity revisited with artifact evidence. Internal validity: the preregistered shared\_initial\_candidate semantics intentionally isolates verifier detection but prevents any verifier-influenced change in executed candidate distributions; this choice explains why a degenerate result may occur in low-error baselines. Construct validity: the emulator-observed fault label is a conservative operational proxy, but emulator-level tests omit traffic-dependent and hardware-specific failure modes; validator traces record final\_state snapshots but cannot capture live-traffic interaction effects. External validity: experiments used a single declared model (gpt-5-mini), a synthetic task bank created by deterministic\_netops\_task\_generator\_v5, and a single deterministic verifier; generalization to other generators, vendor dialects, or production hardware is limited. Conclusion validity: with baseline failure prevalence = 0 the confirmatory paired test lacks the necessary discordance to make inferential claims about verifier usefulness under different baseline regimes.

Practical recommendations grounded in the archive. Based on the archived manifests and traces we recommend that future preregistered studies aiming to quantify verifier utility should:
1) explicitly select or seed a baseline failure prevalence (via adversarial items, ambiguous intents, or historical change examples) so the paired estimand has scope to detect reductions; 2) preregister whether shared\_initial\_candidate semantics or independent-generation semantics will be used, because they answer different operational questions (detection on identical outputs vs. end-to-end improvement under verifier-mediated change); 3) record explicit sampling parameters (non-null temperature) in execution/model\_configuration.json to remove provider-default ambiguity; 4) include repair-enabled arms to measure both detection and net operational improvement; and 5) log verifier-blocking events and, where ethically permissible, execute or simulate blocked candidates to estimate blocking cost. All these recommendations map directly to fields and artifacts present in the current evidence bundle (task manifests, execution/model\_configuration.json, execution/scoring/*.json, and analysis logs) and can be implemented without inventing new tooling.

Concluding perspective. The degenerate confirmatory outcome reported here is an empirically supported datapoint: under the tested constrained, synthetic conditions the generator and manifest constraints produced valid candidates and a canonical verifier did not change executed outcomes. The key scientific value of the run is therefore methodological and procedural: the archived artifacts demonstrate a fully preregistered, deterministic pipeline for measuring verifier detection under tightly controlled semantics, and they make explicit which preregistration choices (shared\_initial\_candidate vs independent generation, sampling parameterization, adversarial seeding, repair permissions) determine whether future confirmatory tests can meaningfully quantify verifier utility.

\section{Limitations}
\begin{itemize}
\item Confirmatory ceiling/null outcome: the baseline generator produced zero emulator faults on the preregistered task set, preventing direct measurement of verifier detection benefit.
\item Shared\_initial\_candidate semantics: the design measures verifier effect on the same generated candidate only and does not assess independent per-condition generation, multi-sample generation, or repairs unless explicitly permitted by the contract.
\item Single-model scope: experiments used one declared model (gpt-5-mini); results do not imply generality across models or versions.
\item Synthetic/emulated tasks: task corpus and emulator executions differ from production devices and live traffic; external validity to production deployments is limited.
\item Sampling-parameter uncertainty: execution/model\_configuration.json shows temperature = null (sampling parameters unset); null indicates provider-default runtime sampling semantics and is a residual uncertainty recorded in the archived configuration.
\item Residual contamination risk: deterministic provenance and exact-match checks were applied, but private training-data contamination of the hosted model cannot be fully ruled out and remains residual uncertainty.
\end{itemize}

\section{Conclusion}
We executed a preregistered paired evaluation (c1-gnr-vv-2026-0001) under shared\_initial\_candidate semantics to test whether a canonical deterministic verifier reduces the proportion of LLM-generated configurations that would execute with operational faults. For the chosen model (gpt-5-mini), verifier (deterministic\_netops\_validator\_v5), and synthetic/emulated task bank, baseline executions produced zero emulator faults and the verifier did not change outcomes (paired difference = 0.0; bootstrap 95\% CI [0.0,0.0]; exact McNemar p = 1.0). This artifact-grounded null/ceiling outcome is internally consistent with the preregistered design: when identical generated candidates produce no executed faults a verifier cannot reduce executed faults under the shared\_initial\_candidate contract. Measuring verifier utility under realistic error regimes requires preregistered contracts that permit independent or multiple generator samples, repair-enabled flows, adversarial task sets, and multi-model comparisons. The archived artifacts (responses, task manifests, validator traces, execution manifest, and analysis logs) provide a reproducible baseline and a template for those follow-on confirmatory designs and power calculations.

\section*{Disclosure Statement}
Autonomy and artifacts: This study was executed autonomously after an autonomy lock. Human role: a Research Director selected the topic family and approved the immutable initial master prompt prior to the autonomy lock; no human manually edited technical manuscript content after the lock. Models and tools: initial generation used gpt-5-mini (declared\_model\_version recorded in execution/model\_configuration.json) orchestrated via the hosted\_netops\_gvr\_v1 adapter; deterministic\_netops\_validator\_v5 (validator\_version = "5.0") served as the canonical verifier; an isolated emulator executed candidate configurations. Preregistration: study c1-gnr-vv-2026-0001 was preregistered and frozen before observing outcomes. Immutable master prompt reference: provenance/master\_prompt.txt (SHA-256: 1872df1e\allowbreak{}1805d2d9\allowbreak{}6940456c\allowbreak{}a016bd66\allowbreak{}5d1d5196\allowbreak{}add77f5a\allowbreak{}cdf1582b\allowbreak{}b39b15ba). Key archived artifacts include execution/raw\_results.jsonl, execution/task\_manifest.jsonl, execution/model\_configuration.json, analysis/paired\_contingency\_table.csv, analysis/condition\_summary.csv, deterministic\_reconciliation.json, and per-episode validator traces under execution/scoring/. The model configuration records temperature = null (provider-default sampling semantics archived in execution/model\_configuration.json). Provider-call audit: hosted\_model\_call\_event\_count = 72. No human manual manuscript editing or content insertion occurred after the autonomy lock.

\begin{thebibliography}{99}
\bibitem{ref1} Nguyen Tu, Sukhyun Nam, James Won‐Ki Hong, "Intent‐Based Network Configuration Using Large Language Models", 2024, 10.1002/nem.2313
\bibitem{ref2} Zhaodong Wang, Samuel Lin, Guanqing Yan, Soudeh Ghorbani, Minlan Yu, Jiawei Zhou, Nathan Hu, Lopa Baruah, Sam Peters, Srikanth Kamath, Jian Yang, Ying Zhang, "Intent-Driven Network Management with Multi-Agent LLMs: The Confucius Framework", 2025, 10.1145/3718958.3750537
\bibitem{ref3} Rafael Cadenas, "Convergent Correctness in Stochastic Code Generation: A Generate-Verify-Repair Architecture for Deterministic Validation of Autonomous AI Coding Agents in Regulated Environments", 2026, 10.2139/ssrn.6754899
\bibitem{ref4} http://citeseerx.ist.psu.edu/viewdoc/summary?doi=10.1.1.300.8659
\end{thebibliography}

\end{document}
